% ============================================================
\chapter{Indexation et Structures de Donn\'ees}
\label{ch:indexation}
% ============================================================

Imaginez une biblioth\`eque de dix millions de livres sans catalogue. Pour trouver un ouvrage, il faudrait parcourir chaque \'etag\`ere, un livre \`a la fois. C'est exactement ce que fait un SGBD lors d'un \emph{full table scan}~: il examine chaque ligne de la table. Un \emph{index} est le catalogue de la base de donn\'ees~: une structure auxiliaire qui permet de localiser les lignes pertinentes en $O(\log n)$ au lieu de $O(n)$. L'arbre B+ (B+ tree), invent\'e par Rudolf Bayer et Edward McCreight en 1972, est la structure d'indexation dominante~: \'equilibr\'e, efficace pour les recherches par \'egalit\'e et par intervalle, et optimis\'e pour les acc\`es disque. Les index hash, les index bitmap et les index spatiaux (R-tree) compl\`etent l'arsenal. Ce chapitre d\'eveloppe ces structures et les strat\'egies d'indexation qui font la diff\'erence entre une requ\^ete en millisecondes et une requ\^ete en minutes.

\begin{intuition}[title=\textbf{Id\'ee centrale}]
Un index est une structure de donn\'ees auxiliaire qui acc\'el\`ere la
recherche dans une table en \'evitant le parcours s\'equentiel
(\emph{full table scan}). Le choix du bon type d'index et des bonnes
colonnes \`a indexer est essentiel pour les performances.
\end{intuition}

% ------------------------------------------------------------
\section{Motivation}
% ------------------------------------------------------------

Sans index, toute recherche dans une table de $n$ lignes n\'ecessite
un parcours s\'equentiel en $O(n)$. Un index bien choisi r\'eduit la
complexit\'e \`a $O(\log n)$ voire $O(1)$.

\begin{remarque}
Un index acc\'el\`ere les lectures mais ralentit les \'ecritures (insertions,
mises \`a jour, suppressions) car la structure doit \^etre maintenue.
Le compromis lecture/\'ecriture guide le choix des index.
\end{remarque}

% ------------------------------------------------------------
\section{Arbres B et B$^+$}
% ------------------------------------------------------------

\begin{definition}[Arbre B d'ordre $m$]
\index{arbre B}
Un \textbf{arbre B} d'ordre $m$ est un arbre de recherche \'equilibr\'e
o\`u chaque n\oe ud interne a au plus $m$ fils et au moins $\lceil m/2 \rceil$
fils (sauf la racine). Les cl\'es sont tri\'ees dans chaque n\oe ud et les
donn\'ees sont stock\'ees dans tous les n\oe uds.
\end{definition}

\begin{definition}[Arbre B$^+$]
\index{arbre B+}
Un \textbf{arbre B$^+$} est une variante o\`u~:
\begin{itemize}
  \item les donn\'ees ne sont stock\'ees que dans les \textbf{feuilles},
  \item les n\oe uds internes ne contiennent que des cl\'es de routage,
  \item les feuilles sont cha\^in\'ees en liste doublement cha\^in\'ee
        pour les parcours s\'equentiels.
\end{itemize}
\end{definition}

\begin{proposition}[Complexit\'e de l'arbre B$^+$]
Pour un arbre B$^+$ d'ordre $m$ contenant $n$ cl\'es~:
\begin{itemize}
  \item Hauteur~: $h = O(\log_m n)$.
  \item Recherche, insertion, suppression~: $O(\log_m n)$ acc\`es disque.
  \item Parcours s\'equentiel d'un intervalle de $k$ cl\'es~: $O(\log_m n + k/B)$
        o\`u $B$ est le nombre de cl\'es par feuille.
\end{itemize}
\end{proposition}

\begin{attention}
L'arbre B$^+$ est la structure d'indexation par d\'efaut de la plupart
des SGBD relationnels (PostgreSQL, MySQL/InnoDB, Oracle, SQL Server).
Ne pas confondre avec l'arbre binaire de recherche (BST) qui n'est pas
adapt\'e aux acc\`es disque.
\end{attention}

\begin{minted}{python}
-- Creation d'un index B+ (implicite) en SQL
CREATE INDEX idx_nom ON etudiants (nom);

-- Index sur plusieurs colonnes (index composite)
CREATE INDEX idx_nom_prenom ON etudiants (nom, prenom);

-- Verifier l'utilisation de l'index
EXPLAIN ANALYZE SELECT * FROM etudiants WHERE nom = 'Dupont';
\end{minted}

% ------------------------------------------------------------
\section{Index de hachage}
% ------------------------------------------------------------

\begin{definition}[Index de hachage]
\index{index de hachage}
Un \textbf{index de hachage} utilise une fonction $h : \text{cl\'e} \to
\{0, 1, \ldots, B-1\}$ pour mapper directement une cl\'e \`a un compartiment
(\emph{bucket}). La recherche par \'egalit\'e est en $O(1)$ en moyenne.
\end{definition}

\begin{remarque}
Les index de hachage ne supportent \textbf{pas} les requ\^etes de plage
(\texttt{WHERE x BETWEEN a AND b}) ni le tri. Ils sont adapt\'es uniquement
aux recherches par \'egalit\'e exacte.
\end{remarque}

\begin{definition}[Hachage extensible]
\index{hachage extensible}
Le \textbf{hachage extensible} (\emph{extendible hashing}) utilise un
r\'epertoire de pointeurs index\'es par les $d$ premiers bits du hachage.
Lorsqu'un compartiment d\'eborde, on double le r\'epertoire (ou on divise
le compartiment) sans r\'eorganiser toute la table.
\end{definition}

% ------------------------------------------------------------
\section{Index bitmap}
% ------------------------------------------------------------

\begin{definition}[Index bitmap]
\index{index bitmap}
Un \textbf{index bitmap} associe \`a chaque valeur distincte $v$ d'une
colonne un vecteur de bits $B_v$ de longueur $n$ (nombre de lignes)~:
$B_v[i] = 1$ si la ligne $i$ a la valeur $v$, 0 sinon.
\end{definition}

\begin{proposition}[Efficacit\'e des op\'erations logiques]
Les requ\^etes combinant plusieurs pr\'edicats se traduisent en
op\'erations bit \`a bit rapides~:
\begin{itemize}
  \item \texttt{WHERE couleur = 'rouge' AND taille = 'L'} $\to$ $B_{\text{rouge}} \wedge B_{\text{L}}$.
  \item \texttt{WHERE couleur = 'rouge' OR taille = 'L'} $\to$ $B_{\text{rouge}} \vee B_{\text{L}}$.
\end{itemize}
Le comptage se fait par \texttt{POPCOUNT} en $O(n/w)$ o\`u $w$ est la taille
du mot machine.
\end{proposition}

\begin{remarque}
Les index bitmap sont particuli\`erement efficaces pour les colonnes \`a
\textbf{faible cardinalit\'e} (sexe, statut, cat\'egorie). Ils sont
utilis\'es massivement dans les entrep\^ots de donn\'ees (Oracle, Vertica).
\end{remarque}

% ------------------------------------------------------------
\section{Index spatiaux~: R-tree}
% ------------------------------------------------------------

\begin{definition}[R-tree]
\index{R-tree}
Un \textbf{R-tree} est un arbre \'equilibr\'e pour l'indexation de donn\'ees
spatiales multidimensionnelles. Chaque n\oe ud contient un \textbf{rectangle
englobant minimum} (MBR) qui englobe tous les objets de son sous-arbre.
\end{definition}

\begin{exemple}
Recherche spatiale~: \og Trouver tous les restaurants dans un rayon de
500m \fg~:
\begin{minted}{python}
-- PostGIS (extension spatiale de PostgreSQL)
CREATE INDEX idx_geom ON restaurants USING GIST (geom);

SELECT nom, ST_Distance(geom, ST_MakePoint(2.35, 48.86)::geography) AS dist
FROM restaurants
WHERE ST_DWithin(geom, ST_MakePoint(2.35, 48.86)::geography, 500)
ORDER BY dist;
\end{minted}
\end{exemple}

% ------------------------------------------------------------
\section{Index plein texte}
% ------------------------------------------------------------

\begin{definition}[Index invers\'e]
\index{index invers\'e}
Un \textbf{index invers\'e} (\emph{inverted index}) associe \`a chaque
terme (mot) la liste des documents (lignes) qui le contiennent, avec
\'eventuellement la position et la fr\'equence.
\end{definition}

\begin{exemple}
\begin{minted}{python}
-- PostgreSQL : index GIN pour la recherche plein texte
CREATE INDEX idx_fts ON articles USING GIN (to_tsvector('french', contenu));

SELECT titre
FROM articles
WHERE to_tsvector('french', contenu) @@ to_tsquery('french', 'bayesien & statistique');
\end{minted}
\end{exemple}

% ------------------------------------------------------------
\section{Strat\'egies de s\'election d'index}
% ------------------------------------------------------------

\begin{definition}[Probl\`eme de s\'election d'index]
\index{s\'election d'index}
\'Etant donn\'e une charge de travail (ensemble de requ\^etes avec leurs
fr\'equences) et un budget d'espace disque, choisir l'ensemble d'index
qui minimise le co\^ut total d'ex\'ecution.
\end{definition}

R\`egles heuristiques~:
\begin{enumerate}
  \item Indexer les colonnes apparaissant dans les clauses \texttt{WHERE},
        \texttt{JOIN} et \texttt{ORDER BY}.
  \item Pr\'ef\'erer les index composites couvrants (\emph{covering index})
        pour \'eviter l'acc\`es \`a la table.
  \item \'Eviter d'indexer les colonnes \`a tr\`es haute cardinalit\'e
        modifi\'ees fr\'equemment.
  \item Utiliser \texttt{EXPLAIN ANALYZE} pour v\'erifier l'utilisation
        effective des index.
\end{enumerate}

\begin{exemple}
\begin{minted}{python}
-- Index couvrant : toutes les colonnes de la requete sont dans l'index
CREATE INDEX idx_covering ON commandes (client_id, date_commande)
    INCLUDE (montant);

-- La requete suivante peut etre satisfaite sans acceder a la table
SELECT date_commande, montant
FROM commandes
WHERE client_id = 42
ORDER BY date_commande DESC;
\end{minted}
\end{exemple}

% ------------------------------------------------------------
\section{Formules cl\'es}
% ------------------------------------------------------------

\begin{formules}
\begin{itemize}
  \item \textbf{B$^+$-tree}~: recherche en $O(\log_m n)$ acc\`es disque, parcours de plage en $O(\log_m n + k/B)$
  \item \textbf{Hash}~: recherche par \'egalit\'e en $O(1)$, pas de requ\^etes de plage
  \item \textbf{Bitmap}~: op\'erations logiques en $O(n/w)$, id\'eal pour faible cardinalit\'e
  \item \textbf{R-tree}~: recherche spatiale en $O(\log n)$ en moyenne
  \item \textbf{Compromis}~: acc\'el\'eration lecture vs.\ surco\^ut \'ecriture
\end{itemize}
\end{formules}

% ------------------------------------------------------------
\section{Exercices}
% ------------------------------------------------------------

\begin{exercice}
Dessiner l'arbre B$^+$ d'ordre 4 obtenu apr\`es insertion s\'equentielle
des cl\'es $3, 7, 1, 14, 8, 5, 11, 17, 13, 6, 23, 12, 20, 26, 4, 16$.
Montrer l'\'etat apr\`es la suppression de la cl\'e 14.
\end{exercice}

\begin{exercice}
Une table \texttt{employes} a 10 millions de lignes. La colonne
\texttt{departement} a 50 valeurs distinctes, la colonne \texttt{salaire}
est continue. Quel type d'index recommandez-vous pour chacune~? Justifier.
\end{exercice}

\begin{exercice}
\'Ecrire une requ\^ete SQL qui ne peut \emph{pas} utiliser un index de
hachage mais peut utiliser un index B$^+$. Expliquer pourquoi.
\end{exercice}

\begin{exercice}
On dispose d'un index bitmap sur la colonne \texttt{statut} (3 valeurs~:
\texttt{actif}, \texttt{inactif}, \texttt{suspendu}) et d'un index bitmap
sur \texttt{region} (10 valeurs). Combien d'op\'erations \texttt{AND} bit
\`a bit sont n\'ecessaires pour r\'epondre \`a \texttt{WHERE statut = 'actif'
AND region IN ('IDF', 'PACA')}~?
\end{exercice}

\begin{exercice}
Utiliser \texttt{EXPLAIN ANALYZE} dans PostgreSQL pour comparer les
performances d'une requ\^ete avec et sans index. Mesurer le nombre de
pages lues (\emph{buffers}) et le temps d'ex\'ecution.
\end{exercice}
